\chapter{Marco teórico}
\label{chap:marco_teorico}

\section{Orígenes} El fantasy es un juego que tiene más años de los que se cree de forma popular, ya después de la segunda guerra mundial empezó la primera forma, más rudimentaria y primitiva de juego fantasy, que fué evolucionando, añadiendo estadísticas y sistemas de puntuación hasta transformarse en lo que es hoy en día.

    \subsection{El nacimiento del concepto Fantasy: La GOPPPL y la Rotisserie League} 
    El primer juego fantasy lo inició un empresario de Oakland y socio minoritario de los Oakland Riders, Wilfred "Bill" Winkenbach, aunque cualquier parecido con los juegos actuales era una coincidencia, pues no tenían nada que ver con las complejas webs y aplicaciones actuales. Este inició en la segunda mitad de la década de los 50' el deporte donde inició este tipo de juegos fue el golf, los jugadores seleccionaban un grupo de golfistas y ganaba aquel que consiguiese el equipo con menos golpes, una fórmula sencilla con la que comenzó este fenómeno, posteriormente en 1962 Bill creó la Greater Oakland Professional Pigskin Pronosticators League (GOPPPL) \cite{origen_gopppl}. 
    
    Posteriormente nacería el fantasy tal y como lo conocemos, con estadísticas avanzadas, ya que fue la primera vez que se seleccionaban deportistas y se seguían sus estadísticas durante una temporada, surgió en un restaurante francés de Nueva York llamado La Rotisserie Française y su creador fue Daniel Okrent, escritor y editor. En 1980 se publicó un artículo en el New York Times sobre la liga, lo que provocó que se propagase por todo Estados Unidos \cite{rotisserie_nyt}, la gente copiaba las reglas y empezaban a jugar.

    \subsection{El primer Fantasy de fútbol: Fantacalcio} 
    En 1988, Riccardo Albini, un periodista italiano, viajó a Estados Unidos y compró un libro sobre las reglas del Fantasy Football y la Rotisserie League, al volver a Italia, se quedó con la duda de por qué no exisitía eso pero con la Serie A, por lo que estuvo un par de años modificando y adaptando las reglas para el fútbol y finalmente en 1990 publicó el manual del Fantacalcio, el primer juego fantasy de fútbol, para que posteriormente se publicase en la Gazzeta dello Sport el juego oficial \cite{fantacalcio_historia}.

    \subsection{La llegada a España: La Liga Fantástica Marca} El inicio de la locura fantasy en España inició en 1994 cuando el periódico Marca observó en el Fantacalcio una gran iniciativa que todavía no había aterrizado en España, por lo que para la temporada 1994/95 surgió La Liga Fantástica Marca, de esta forma podrían fidelizar a los lectores para que lo compraran todos los días y no solo los lunes para ver los resultados.

    Te daban un presupuesto y tenias que formar un once para la jornada, que se rellenaba con recortes de los cupones que traía el periódico. Esos cupones se enviaban por fax o correo antes del inicio de la jornada, por lo que, si querías cambiar algún lesionado del once de última hora, tenías que llamar a un número de pago para que te hiciesen el cambio.

    \subsection{Actualidad en los juegos Fantasy} En la década de los 2000, comenzó la masificación de Internet lo que conllevó a la evolución de los juegos Fantasy, en lo que ahora conocemos, webs o aplicaciones interactivas, el precursor de esto fué Comunio, creado en Alemania en el año 2000, alcanzó un éxito masivo en España a partir del 2006 \cite{comunio_origen}.

    Ahora los managers no tenían que rellenar cupones del periódico, simplemente entraban a la web y podían crear ligas privadas con amigos y pujar en mercados que se actualizaban diariamente, actualizando así los precios de los jugadores, dependiendo de cuantos managers pujasen y de cuanto pujasen. Estas webs usaban el sistema de picas tradicional ya que por ese entonces, todavía no existían proveedores de datos deportivos al nivel que existen actualmente.

    Actualmente, con el avance de la tecnología y de la recopilación de datos, la mayoría de los juegos fantasy implementan sistemas de puntuación basados en estadísticas o sistemas híbridos.

\section{Sistemas de puntuación}
Los juegos fantasy de fútbol tienen diversos sistemas de puntuación que han ido cambiando a lo largo de los años, fruto también del número de estadísticas de las que se disponían. Debido a que hace más de dos décadas no era común que se recogiesen datos como los balones perdidos, las faltas cometidas o los tiros a puerta, los sistemas de puntuación no dependían de ellos sino de cronistas, que eran asignados a los partidos por los diferentes periódicos o webs que gestionaban el fantasy, por lo que podemos dividir los sistemas en varios grupos.

\begin{itemize}
    \item Los sistemas de puntuación basados en estadísticas son aquellos cuyo sistema da puntos en función de la puntuación Sofascore, que se otorga a cada jugador dependiendo de estadísticas como tiros totales, tiros a puerta, pases precisos, pases clave, regates, faltas recibidas, posesiones perdidas, entradas, intercepciones, despejes, etc.

    Algunos ejemplos de este sistema son:
    \begin{itemize}
        \item LaLiga Fantasy
        \item Biwenger Sofascore.
        \item Mister Sofascore.
    \end{itemize}

    \item Los sistemas de puntuación basados en cronistas son aquellos en los que la mayor parte de la puntuación se otorga con un baremo, que traduce de picas a puntos, al que se le añaden puntos extra por eventos mayores como goles, asistencias o tarjetas recibidas.

    Algunos ejemplos de este sistema son:
    \begin{itemize}
        \item Biwenger Cronistas.
        \item Comunio.
        \item Futmondo prensa.
    \end{itemize}

    \item También existen modelos de puntuación mixtos que combinan las métricas avanzadas de los sistemas de estadísticas con una parte de puntuación dada por los cronistas. Generalmente estos sistemas añaden un apartado que da entre 0 y 4 puntos basado la puntuación que le da el cronista elegido para el partido al partido del jugador.

    Algunos ejemplos de este sistema son:
    \begin{itemize}
        \item Fantasy Marca.
    \end{itemize}
\end{itemize}

Para poder comprender de mejor forma los distintos sistemas de puntuación, a continuación se podrán ver tablas con las equivalencias de los sistemas a puntos fantasy (estas tablas están basadas en plataformas líderes como Biwenger) \cite{reglas_biwenger}. Tras analizar estas tablas se podrá, determinar el mejor modelo para el dataset usado en este proyecto, el cuál se expondrá y analizará posteriormente en este documento.

\subsubsection{Comparativa de Baremos de Puntuación}

El modelo basado en cronistas (Tabla \ref{tab:puntos_picas}) destaca por su simplicidad y su subjetividad que dependiendo del contexto del partido traduce de mejor manera la actuación del jugador en el partido a puntos.

\begin{table}[H]
    \centering
    \begin{tabular}{|l|l|c|}
        \hline
        \textbf{Valoración del Cronista} & \textbf{Significado} & \textbf{Puntos Fantasy} \\
        \hline
        Sin calificar (S.C.) & Tiempo insuficiente en el campo & 0 \\
        \hline
        0 Picas (-) & Mal partido o expulsión temprana & -2 \\
        \hline
        1 Pica ($\spadesuit$) & Partido regular o aceptable & 2 \\
        \hline
        2 Picas ($\spadesuit\spadesuit$) & Buen partido & 6 \\
        \hline
        3 Picas ($\spadesuit\spadesuit\spadesuit$) & Partido excepcional & 10 \\
        \hline
        4 Picas ($\spadesuit\spadesuit\spadesuit\spadesuit$) & Actuación histórica o perfecta & 14 \\
        \hline
    \end{tabular}
    \caption{Sistema de puntuación tradicional basado en cronistas deportivos (Picas).}
    \label{tab:puntos_picas}
\end{table}

Por el contrario, el sistema estadístico (Tabla \ref{tab:puntos_sofascore}) depende de un motor de \textit{Big Data} (como Sofascore u Opta) que evalúa decenas de métricas en tiempo real (pases clave, duelos ganados, pérdidas, faltas recibidas, toques, regates, despejes, etc). Este motor calcula una nota decimal del 1 al 10, que posteriormente el juego \textit{Fantasy} traduce a puntos, ofreciendo una precisión mucho mayor en términos generales de la actuación del jugador en el partido.

\begin{table}[H]
    \centering
    \begin{tabular}{|c|c|c|c|}
        \hline
        \textbf{Nota Estadística} & \textbf{Puntos} & \textbf{Nota Estadística} & \textbf{Puntos} \\
        \hline
        0 -- 4.9 & -6 & 7.0 -- 7.1 & 5 \\ \hline
        5.0 -- 5.1 & -5 & 7.2 -- 7.3 & 6 \\ \hline
        5.2 -- 5.3 & -4 & 7.4 -- 7.5 & 7 \\ \hline
        5.4 -- 5.5 & -3 & 7.6 -- 7.7 & 8 \\ \hline
        5.6 -- 5.7 & -2 & 7.8 -- 7.9 & 9 \\ \hline
        5.8 -- 5.9 & -1 & 8.0 -- 8.1 & 10 \\ \hline
        6.0 -- 6.1 & 0 & 8.2 -- 8.5 & 11 \\ \hline
        6.2 -- 6.3 & 1 & 8.6 -- 8.9 & 12 \\ \hline
        6.4 -- 6.5 & 2 & 9.0 -- 9.4 & 13 \\ \hline
        6.6 -- 6.7 & 3 & 9.5 -- 10 & 14 \\ \hline
        6.8 -- 6.9 & 4 & \multicolumn{2}{c|}{} \\ \hline
    \end{tabular}
    \caption{Mapeo oficial de nota estadística SofaScore a puntos Fantasy en Biwenger.}
    \label{tab:puntos_sofascore}
\end{table}

Adicionalmente al sistema de puntuación base elegido, cronistas o estadísticas, todos los sistemas implementan una capa extra de \textbf{Eventos Fijos} (Tabla \ref{tab:eventos_fijos}). Estos eventos suman o restan puntos basándose en acontecimientos verificables que ocurren en los partidos, cuyo valor varía en función de la posición que le es asignada por el juego.

\begin{table}[H]
    \centering
    \begin{tabular}{|l|c|c|c|c|}
        \hline
        \textbf{Evento de Partido} & \textbf{Portero} & \textbf{Defensa} & \textbf{Medio} & \textbf{Delantero} \\
        \hline
        Gol anotado & +6 & +5 & +4 & +3 \\
        \hline
        Asistencia de gol & +1 & +1 & +1 & +1 \\
        \hline
        Penalti fallado & -2 & -2 & -2 & -2 \\
        \hline
        Penalti parado & +3 & +3 & +3 & +3 \\
        \hline
        Tarjeta Amarilla & -1 & -1 & -1 & -1 \\
        \hline
        Tarjeta Roja & -3 & -3 & -3 & -3 \\
        \hline
    \end{tabular}
    \caption{Bonificaciones y penalizaciones por eventos fijos según la posición del jugador.}
    \label{tab:eventos_fijos}
\end{table}

Además, algunas plataformas modernas oficiales, como \textbf{LaLiga Fantasy}, el cuál es el juego fantasy oficial de LaLiga, optan por no hacer uso de una nota decimal global y prefieren un modelo de micro-eventos \cite{reglas_laligafantasy}. En este sistema, el algoritmo suma y resta puntos constantemente basándose en métricas avanzadas de \textit{Big Data} (Tabla \ref{tab:puntuacion_laliga_fantasy}).

\begin{table}[H]
    \centering
    \small
    \begin{tabular}{|p{4.5cm}|c|c|c|c|}
        \hline
        \textbf{Acción Evaluada} & \textbf{PT} & \textbf{DF} & \textbf{MC} & \textbf{DL} \\
        \hline
        \textbf{Minutos Jugados} & \multicolumn{4}{c|}{} \\ \hline
        Menos de 60 minutos & +1 & +1 & +1 & +1 \\ \hline
        Más de 60 minutos & +2 & +2 & +2 & +2 \\ \hline
        \textbf{Goles y Asistencias} & \multicolumn{4}{c|}{} \\ \hline
        Gol anotado & +6 & +6 & +5 & +4 \\ \hline
        Asistencia de gol & +3 & +3 & +3 & +3 \\ \hline
        Ocasión manifiesta (Pase clave) & +1 & +1 & +1 & +1 \\ \hline
        \textbf{Defensa y Portería} & \multicolumn{4}{c|}{} \\ \hline
        Portería a cero ($>$ 60 min) & +4 & +3 & +2 & +1 \\ \hline
        Goles recibidos (Cada 2 goles) & -2 & -2 & -1 & -1 \\ \hline
        Paradas (Cada 2 paradas) & +1 & - & - & - \\ \hline
        Balones recuperados (Cada 5) & +1 & +1 & +1 & +1 \\ \hline
        Despejes (Cada 3) & +1 & +1 & +1 & +1 \\ \hline
        \textbf{Penaltis y Sanciones} & \multicolumn{4}{c|}{} \\ \hline
        Penalti parado & +5 & - & - & - \\ \hline
        Penalti provocado a favor & +2 & +2 & +2 & +2 \\ \hline
        Penalti cometido o fallado & -2 & -2 & -2 & -2 \\ \hline
        Tarjeta Amarilla & -1 & -1 & -1 & -1 \\ \hline
        Tarjeta Roja & -3 & -3 & -3 & -3 \\ \hline
        \textbf{Bonus de Ataque} & \multicolumn{4}{c|}{} \\ \hline
        Tiros a puerta (Cada 2) & +1 & +1 & +1 & +1 \\ \hline
        Regates logrados (Cada 2) & +1 & +1 & +1 & +1 \\ \hline
        Balones al área (Cada 2) & +1 & +1 & +1 & +1 \\ \hline
        \textbf{Penalizaciones de Posesión} & \multicolumn{4}{c|}{} \\ \hline
        Pérdidas de balón (-1 punto) & Cada 8 & Cada 8 & Cada 10 & Cada 12 \\ \hline
        \textbf{Nota Subjetiva Adicional} & \multicolumn{4}{c|}{} \\ \hline
        Puntos DAZN (Impacto en partido) & 0 a +4 & 0 a +4 & 0 a +4 & 0 a +4 \\ \hline
    \end{tabular}
    \caption{Sistema oficial de puntuación de LaLiga Fantasy detallado por demarcación (PT: Portero, DF: Defensa, MC: Medio, DL: Delantero).}
    \label{tab:puntuacion_laliga_fantasy}
\end{table}

La inclusión de estas métricas anteriormente mencionadas (Tabla \ref{tab:puntuacion_laliga_fantasy}) demuestra la evolución de la granularidad de los datos recogidos por las empresas encargadas de conseguirlos y tratarlos. Mientras que un juego 'antiguo' puede funcionar leyendo un simple resumen de partido en el que aparecen datos como los goles, asistencias y tarjetas, un juego 'moderno' en muchas ocasiones requiere de procesar archivos o bases de datos que registren cada acción individual del jugador sobre el terreno de juego.

\section{La Inteligencia Artificial en el Ecosistema del Desarrollo de Software}
\label{sec:ia_desarrollo_software}

En los últimos años, el proceso de desarrollo de software ha experimentado una transformación profunda debido a la irrupción de la Inteligencia Artificial Generativa y los Modelos de Lenguaje de Gran Tamaño, LLMs, por sus siglas en inglés. Desde la aparición de los modelos de lenguaje capaces de aprender a partir de unos pocos ejemplos \cite{brown_fewshot}, esta tecnología ha dejado de ser una mera herramienta experimental para convertirse en un componente completamente fundamental en los flujos de trabajo de ingeniería de software modernos.

\subsection{Estado Actual de la Inteligencia Artificial Generativa}
Los modelos actuales, tales como los pertenecientes a las familias GPT de OpenAI, Claude de Anthropic y Gemini de Google, demuestran capacidades avanzadas de razonamiento lógico, traducción de lenguajes y generación de código estructurado. La capacidad de estos modelos para comprender el contexto de bases de código completas, analizar dependencias y sugerir algoritmos complejos ha propiciado su integración directa en entornos de desarrollo integrado, IDEs, a través de herramientas de asistencia como Antigravity, Cursor o Claude Code.

\subsection{Uso de la IA en Proyectos de Ingeniería de Software}
En los procesos de aquellos proyectos que implican el manejo de grandes volúmenes de datos e interacciones en tiempo real, como \textit{Retro Fantasy}, el uso de asistentes de IA se orienta a optimizar varias fases críticas de este proceso:
\begin{itemize}
    \item \textbf{Generación y optimización de código de bajo nivel}: Automatización de tareas repetitivas o código de infraestructura (controladores, configuraciones de API y modelos de datos) permitiendo al desarrollador enfocarse en la arquitectura general y la lógica de negocio.
    \item \textbf{Diseño de consultas de base de datos y optimización}: La IA facilita la creación y verificación de scripts SQL complejos, políticas de seguridad a nivel de fila (RLS) en sistemas relacionales como PostgreSQL y el diseño de disparadores (\textit{triggers}).
    \item \textbf{Auditoría y refactorización}: Ayuda en la identificación de fallos de lógica, redundancia de código y vulnerabilidades, permitiendo así al desarrollador aplicar principios de diseño como \textit{SOLID} y \textit{Clean Architecture} además de buenas prácticas de mejor forma.
\end{itemize}

Esta sinergia entre el desarrollador humano (que actúa como arquitecto y supervisor de la lógica del sistema) y el asistente de Inteligencia Artificial (como redactor técnico de código) define un nuevo paradigma de programación en pareja, \textit{Pair Programming}, que acorta los tiempos de desarrollo y aumenta la fiabilidad y el alcance del producto final. Estudios empíricos sobre asistentes de código como GitHub Copilot constatan mejoras significativas en la productividad de los desarrolladores que adoptan este tipo de herramientas \cite{peng_copilot}.